\documentclass[11pt]{article}

\usepackage[a4paper,margin=1in]{geometry}
\usepackage{amsmath,amssymb,amsthm,mathtools,bm}
\usepackage{enumitem}
\usepackage{booktabs,tabularx,array}
\usepackage{microtype}
\usepackage[hidelinks]{hyperref}
\usepackage{graphicx}
\usepackage[section]{placeins}
\usepackage{xcolor}
\usepackage{tikz}
\usetikzlibrary{arrows.meta,positioning,shapes.geometric}

\newtheorem{assumption}{Assumption}[section]
\newtheorem{definition}[assumption]{Definition}
\newtheorem{proposition}[assumption]{Proposition}
\newtheorem{theorem}[assumption]{Theorem}
\newtheorem{lemma}[assumption]{Lemma}
\newtheorem{corollary}[assumption]{Corollary}
\theoremstyle{remark}
\newtheorem{remark}[assumption]{Remark}
\newtheorem{example}[assumption]{Example}

\newcommand{\E}{\mathbb E}
\newcommand{\Pp}{\mathbb P}
\newcommand{\R}{\mathbb R}
\newcommand{\N}{\mathbb N}
\newcommand{\Law}{\mathcal L}
\newcommand{\cH}{\mathcal H}
\newcommand{\cW}{\mathcal W}
\newcommand{\cT}{\mathcal T}
\newcommand{\cL}{\mathcal L}
\newcommand{\cC}{\mathcal C}
\newcommand{\norm}[1]{\left\lVert #1\right\rVert}
\newcommand{\abs}[1]{\left\lvert #1\right\rvert}
\newcommand{\ip}[2]{\left\langle #1,#2\right\rangle}
\newcommand{\op}{\mathrm{o}_{\Pp}}
\newcommand{\Op}{\mathrm{O}_{\Pp}}
\newcommand{\cf}{\mathrm{cf}}
\newcommand{\fact}{\mathrm{fact}}
\newcommand{\Ph}{\Phi}
\newcommand{\ph}{\phi}
\newcommand{\Wtwo}{\mathsf W_2}
\newcommand{\Lam}{\Lambda}
\newcommand{\an}{a_{N,n}}
\newcommand{\bn}{b_{N,n}}
\newcommand{\LamN}{\Lambda_N}
\newcommand{\thetahat}{\widehat\theta_{N,n}}
\newcommand{\sigmahat}{\widehat\sigma_{N,n}}

\title{Work Package 7: Galerkin, Temporal, Operator, and Total Inferential Error\\
\large Explicit field-functional rates, numerical-bias regimes, and root-information-valid inference}
\author{}
\date{7 July 2026}

\begin{document}
\maketitle

\begin{abstract}
This note closes the positive theorem chain for causal operator interventions in linear parabolic SPDEs. The exact Galerkin semantics of Work Package 1, the mixed-rate likelihood theory of Work Package 3, the uniform target processes and multiplier bands of Work Packages 4--5, and the operator response calculus of Work Package 6 are combined into a single approximation and inference theorem. For the compact space-time weight class used in the uniform process theory, explicit deterministic-mean and stochastic-variance tails are derived. A structural distinction emerges: the strong shared-noise field-functional error contains the square root of the variance tail, whereas non-degenerate one-dimensional Gaussian laws admit sharper optimal $W_2$, Kolmogorov, and quantile rates because the difference of standard deviations is linear in the variance tail. Temporal approximation is treated in two forms: an exact-modal functional scheme based on projection of the deterministic weight, and a left-point/state-discretised scheme under temporal regularity. Bounded operator-approximation error is propagated through the Work Package 6 derivative, with a second-order rate after first-order response correction. The resulting target error is the sum of statistical, spatial, temporal, and operator terms. We identify undersmoothing, local numerical bias, and dominant bias regimes; derive bias-enlarged simultaneous bands; and give cost-optimal allocations of the spatial level, time step, and operator resolution. Numerical experiments recover the predicted spatial slopes $-2.6$, $-4.6$, and $-2.3$, a first-order temporal law rate, and the coverage consequences of balanced versus undersmoothed approximation. This package passes Decision Gate 6 and authorises full-field inference for the positive Gaussian core.
\end{abstract}

\section{Programme interface and model}

The paper's positive theorem chain is now
\[
\begin{aligned}
 &\text{semantics}\Rightarrow\text{identification}\Rightarrow\text{LAN and efficiency}\Rightarrow\text{uniform target processes}\\
 &\hspace{18mm}\Rightarrow\text{feasible bands}\Rightarrow\text{total numerical--statistical error}.
\end{aligned}
\]
The present work package proves the final implication.

Let $H=L^2(D)$ and set $\Lam=-A_0$. Assume
\begin{equation}\label{eq:spectrum}
 \Lam e_k=\lambda_ke_k,
 \qquad
 0<\lambda_1\le\lambda_2\le\cdots,
 \qquad
 \lambda_k\asymp k^{2/d}.
\end{equation}
For an absolute intervention $a\in\widetilde\Theta=[a_-,a_+]\Subset(0,\infty)$,
\begin{equation}\label{eq:spde}
 dU_t^a=-a\Lam U_t^a\,dt+\sigma Q^{1/2}dW_t,
 \qquad U_0^a=u_0,
\end{equation}
where $Q^{1/2}e_k=q_ke_k$ and $\sigma\in[\sigma_-,\sigma_+]\Subset(0,\infty)$.

For $\cH=L^2(0,T;H)$ and $w\in\cH$, define
\begin{equation}\label{eq:X}
 X_{a,w}:=\Psi_w(U^a)
 =\int_0^T\ip{w_t}{U_t^a}\,dt.
\end{equation}
It is Gaussian:
\begin{equation}\label{eq:law}
 X_{a,w}\sim
 \mathcal N\big(m_a(w),\sigma^2V_a(w)\big),
\end{equation}
with
\begin{align}
 m_a(w)
 &=\sum_{k\ge1}u_{0,k}\int_0^T w_k(t)e^{-a\lambda_kt}\,dt,
 \label{eq:mean}\\
 V_a(w)
 &=\sum_{k\ge1}q_k^2\int_0^T
 \left[
 \int_s^T w_k(t)e^{-a\lambda_k(t-s)}\,dt
 \right]^2ds.
 \label{eq:variance}
\end{align}

\begin{assumption}[Target and regularity class]\label{ass:class}
There exist $r,s>0$, $R<\infty$, $r_0\ge0$, and $\beta\in\R$ such that
\begin{align}
 \cW_R^{r,s}
 &:={\left\{
 w=\sum_{k,\ell}w_{k\ell}\varphi_\ell e_k:
 \sum_{k,\ell}
 (1+\lambda_k)^{2s}(1+\nu_\ell)^{2r}
 |w_{k\ell}|^2\le R^2
 \right\}},
 \label{eq:ellipsoid}\\
 u_0&\in D(\Lam^{r_0/2}),\label{eq:u0}\\
 q_k^2&\le C_q\lambda_k^{-\beta},
 \qquad \beta+1>\frac d2,\label{eq:qdecay}\\
 \sum_{k\ge1}\frac{q_k^2}{\lambda_k}&<\infty.\label{eq:baseline-noise}
\end{align}
The operative class $\cW\subset\cW_R^{r,s}$ satisfies
\begin{equation}\label{eq:variance-lower}
 0<v_-\le V_a(w)\le v_+<\infty
 \qquad
 \text{for all }(a,w)\in\widetilde\Theta\times\cW.
\end{equation}
\end{assumption}

The temporal weight in \eqref{eq:ellipsoid} is needed for compactness and temporal approximation. The lower variance bound is inherited from Work Package 4 and is the mechanism that upgrades the stochastic variance tail into a sharper Gaussian-law rate.

\section{Exact Galerkin primitives}

Let $P_K$ be the spectral projection onto $\operatorname{span}\{e_1,\dots,e_K\}$ and define
\begin{equation}\label{eq:XK}
 X_{a,w}^{K}:=\Psi_{P_Kw}(P_KU^a).
\end{equation}
Its mean and primitive variance are
\begin{align}
 m_{a,K}(w)
 &=\sum_{k\le K}u_{0,k}\int_0^T w_k(t)e^{-a\lambda_kt}\,dt,
 \label{eq:mK}\\
 V_{a,K}(w)
 &=\sum_{k\le K}q_k^2\int_0^T
 \left[
 \int_s^T w_k(t)e^{-a\lambda_k(t-s)}\,dt
 \right]^2ds.
 \label{eq:VK}
\end{align}
In particular, $V_a(w)-V_{a,K}(w)\ge0$.

\begin{theorem}[Uniform deterministic mean tail]\label{thm:mean-tail}
Under Assumption~\ref{ass:class}, with
\begin{equation}\label{eq:mum}
 \mu_m:=\frac{r_0+1+2s}{2},
\end{equation}
there exists $C<\infty$ such that
\begin{equation}\label{eq:mean-tail}
 \sup_{a\in\widetilde\Theta,\,w\in\cW}
 |m_a(w)-m_{a,K}(w)|
 \le C\lambda_{K+1}^{-\mu_m}.
\end{equation}
\end{theorem}

\begin{proof}
Since
\[
 \norm{e^{-a\lambda_k\cdot}}_{L^2(0,T)}
 \le(2a_-\lambda_k)^{-1/2},
\]
Cauchy--Schwarz gives
\begin{align*}
 |m_a-m_{a,K}|
 &\le C\sum_{k>K}|u_{0,k}|\lambda_k^{-1/2}
orm{w_k}_{L^2}\\
 &\le C
 \left(\sum_{k>K}\lambda_k^{r_0}|u_{0,k}|^2\right)^{1/2}
 \left(\sum_{k>K}\lambda_k^{-r_0-1}\norm{w_k}_{L^2}^2\right)^{1/2}.
\end{align*}
Insert $\lambda_k^{2s}\lambda_k^{-2s}$ in the second factor and bound the tail by $\lambda_{K+1}^{-(r_0+1+2s)/2}$. Assumption~\ref{ass:class} supplies the uniform constants.
\end{proof}

\begin{theorem}[Uniform stochastic variance tail]\label{thm:variance-tail}
Set
\begin{equation}\label{eq:muV}
 \mu_V:=\beta+2+2s.
\end{equation}
Then
\begin{equation}\label{eq:variance-tail}
 0\le
 \sup_{a,w}
 \{V_a(w)-V_{a,K}(w)\}
 \le C\lambda_{K+1}^{-\mu_V}.
\end{equation}
\end{theorem}

\begin{proof}
For the backward Volterra kernel, Young's inequality yields
\[
 \left\|
 \int_{\cdot}^T w_k(t)e^{-a\lambda_k(t-\cdot)}dt
 \right\|_{L^2(0,T)}
 \le\frac1{a_-\lambda_k}\norm{w_k}_{L^2(0,T)}.
\]
Therefore
\begin{align*}
 V_a-V_{a,K}
 &\le C\sum_{k>K}q_k^2\lambda_k^{-2}\norm{w_k}_{L^2}^2\\
 &\le C\lambda_{K+1}^{-(\beta+2+2s)}
 \sum_{k>K}\lambda_k^{2s}\norm{w_k}_{L^2}^2.
\end{align*}
The ellipsoid bound completes the proof.
\end{proof}

\section{Strong functional error and Gaussian law error}

Because the omitted stochastic modes are centred and orthogonal,
\begin{equation}\label{eq:strong-exact}
 \E|X_{a,w}-X_{a,w}^{K}|^2
 =|m_a-m_{a,K}|^2+\sigma^2(V_a-V_{a,K}).
\end{equation}

\begin{corollary}[Strong coupled functional rate]\label{cor:strong}
Let
\begin{equation}\label{eq:mustr}
 \mu_{\mathrm{str}}
 :=\min\left\{\mu_m,\frac{\mu_V}{2}\right\}.
\end{equation}
Then
\begin{equation}\label{eq:strong-rate}
 \sup_{a,w,\sigma}
 \norm{X_{a,w}-X_{a,w}^{K}}_{L^2(\Omega)}
 \le C\lambda_{K+1}^{-\mu_{\mathrm{str}}}.
\end{equation}
Under Weyl growth,
\begin{equation}\label{eq:rho-str}
 \rho_{\mathrm{str}}
 =\min\left\{
 \frac{r_0+1+2s}{d},
 \frac{\beta+2+2s}{d}
 \right\},
 \qquad
 \norm{X-X^K}_{L^2}=O(K^{-\rho_{\mathrm{str}}}).
\end{equation}
\end{corollary}

The natural shared-noise coupling is not generally optimal at the law level.

\begin{lemma}[Normal-law comparison]\label{lem:normal-comparison}
Let $G_i=\mathcal N(m_i,s_i^2)$ with $s_i\ge s_->0$. Then
\begin{align}
 \Wtwo^2(G_0,G_1)
 &=|m_1-m_0|^2+|s_1-s_0|^2,
 \label{eq:w2-normal}\\
 \sup_z|G_1(-\infty,z]-G_0(-\infty,z]|
 &\le
 \frac{|m_1-m_0|}{s_-\sqrt{2\pi}}
 +\frac{|s_1-s_0|}{s_-\sqrt{2\pi e}}.
 \label{eq:kol-normal}
\end{align}
For $p\in[p_-,p_+]\Subset(0,1)$,
\begin{equation}\label{eq:q-normal}
 |Q_1(p)-Q_0(p)|
 \le|m_1-m_0|+C_{p_-,p_+}|s_1-s_0|.
\end{equation}
\end{lemma}

\begin{proof}
Equation \eqref{eq:w2-normal} follows from the monotone quantile coupling. For the CDF, interpolate first in the mean and then in $\log s$. The derivatives are bounded by $1/(s_-\sqrt{2\pi})$ and $\sup_x|x\ph(x)|=(2\pi e)^{-1/2}$, respectively. Equation \eqref{eq:q-normal} follows from $Q_i(p)=m_i+s_i\Ph^{-1}(p)$.
\end{proof}

\begin{theorem}[Improved Gaussian law rate]\label{thm:law-rate}
Let
\begin{equation}\label{eq:mulaw}
 \mu_{\mathrm{law}}:=\min\{\mu_m,\mu_V\}.
\end{equation}
For all sufficiently large $K$,
\begin{align}
 \sup_{a,w,\sigma}
 \Wtwo\big(\Law(X_{a,w}),\Law(X_{a,w}^{K})\big)
 &\le C\lambda_{K+1}^{-\mu_{\mathrm{law}}},
 \label{eq:w2-rate}\\
 \sup_{a,w,z,\sigma}|F_{a,w,K}(z)-F_{a,w}(z)|
 &\le C\lambda_{K+1}^{-\mu_{\mathrm{law}}},
 \label{eq:cdf-rate}\\
 \sup_{a,w,p,\sigma}|Q_{a,w,K}(p)-Q_{a,w}(p)|
 &\le C\lambda_{K+1}^{-\mu_{\mathrm{law}}}.
 \label{eq:q-rate}
\end{align}
Under Weyl growth, the common law exponent is
\begin{equation}\label{eq:rhoF}
 \rho_F
 =\min\left\{
 \frac{r_0+1+2s}{d},
 \frac{2(\beta+2+2s)}{d}
 \right\}.
\end{equation}
\end{theorem}

\begin{proof}
Assumption~\ref{ass:class} and Theorem~\ref{thm:variance-tail} imply $V_{a,K}(w)\ge v_-/2$ for large $K$. Therefore
\[
 |\sqrt{V_a}-\sqrt{V_{a,K}}|
 =\frac{V_a-V_{a,K}}{\sqrt{V_a}+\sqrt{V_{a,K}}}
 \le C(V_a-V_{a,K}).
\]
Apply Lemma~\ref{lem:normal-comparison} together with Theorems~\ref{thm:mean-tail}--\ref{thm:variance-tail}.
\end{proof}

\begin{remark}[Why the law rate is sharper]
The strong tail contains $(V-V_K)^{1/2}$, while the optimal coupling of two non-degenerate one-dimensional Gaussian laws contains $\sqrt V-\sqrt{V_K}=O(V-V_K)$. The lower variance bound is essential. The improvement disappears near degenerate laws, which is why the complete $L^2$ weight ball cannot support this uniform conclusion.
\end{remark}

\section{Modal-decay corollaries and a refined time-smooth rate}

\begin{corollary}[Pointwise modal-decay rates]\label{cor:pointwise}
Suppose
\[
 |u_{0,k}|\le C\lambda_k^{-r_0/2},
 \qquad
 \norm{w_k}_{L^2}\le C\lambda_k^{-\gamma/2},
 \qquad
 q_k^2\le C\lambda_k^{-\beta}.
\]
If the displayed tails are summable, then
\begin{align}
 |m-m_K|&=O(K^{-\rho_m^{\mathrm{pt}}}),
 &\rho_m^{\mathrm{pt}}
 &=\frac{r_0+\gamma+1-d}{d},
 \label{eq:ptmean}\\
 V-V_K&=O(K^{-\rho_V^{\mathrm{pt}}}),
 &\rho_V^{\mathrm{pt}}
 &=\frac{2(\beta+\gamma+2)-d}{d}.
 \label{eq:ptvar}
\end{align}
Consequently,
\begin{equation}\label{eq:ptrates}
 \rho_{\mathrm{str}}^{\mathrm{pt}}
 =\min\left\{\rho_m^{\mathrm{pt}},\frac{\rho_V^{\mathrm{pt}}}{2}\right\},
 \qquad
 \rho_F^{\mathrm{pt}}
 =\min\left\{\rho_m^{\mathrm{pt}},\rho_V^{\mathrm{pt}}\right\}.
\end{equation}
\end{corollary}

\begin{proof}
Use the Weyl-tail estimate
\[
 \sum_{k>K}\lambda_k^{-p}
 \lesssim K^{1-2p/d},
 \qquad p>d/2,
\]
with $p=(r_0+\gamma+1)/2$ for the mean and $p=\beta+\gamma+2$ for the variance.
\end{proof}

\begin{remark}[Time-smooth modal weights]
The generic mean bound uses only $L^2$ regularity in time and gains $\lambda_k^{-1/2}$. If $w_k(t)=c_k\bar w(t)$ and the Laplace transform of $\bar w$ gains a full factor $\lambda_k^{-1}$, then
\begin{equation}\label{eq:const-rate}
 \rho_m^{\mathrm{smooth}}
 =\frac{r_0+\gamma+2-d}{d}.
\end{equation}
For constant-in-time weights this is the rate verified in Section~\ref{sec:numerical}. It is a design-specific refinement, not a replacement for Theorem~\ref{thm:mean-tail}.
\end{remark}

\section{Temporal approximation}

The target functional, rather than the whole state path, is the primary numerical object. This allows an exact-modal construction that avoids an unnecessary state-discretisation error.

\subsection{Projected-weight exact-modal scheme}

Let $\Pi_\delta$ be a piecewise-polynomial projection on a partition of mesh $\delta$. Define
\begin{equation}\label{eq:XKdelta}
 X_{a,w}^{K,\delta}:=\Psi_{\Pi_\delta w}(P_KU^a).
\end{equation}
In the diagonal additive model, its Gaussian mean and covariance can be evaluated mode by mode; equivalently, integrated OU cell increments can be sampled exactly.

\begin{assumption}[Temporal projection rate]\label{ass:temporal-projection}
For some $\eta>0$,
\begin{equation}\label{eq:projection-rate}
 \sup_{w\in\cW}
\norm{w-\Pi_\delta w}_{\cH}
 \le C\delta^\eta.
\end{equation}
\end{assumption}

\begin{theorem}[Temporal functional and law rate]\label{thm:temporal-projection}
Under Assumptions~\ref{ass:class} and \ref{ass:temporal-projection}, uniformly in $K,a,w,\sigma$,
\begin{align}
 \norm{X_{a,w}^{K}-X_{a,w}^{K,\delta}}_{L^2(\Omega)}
 &\le C\delta^\eta,
 \label{eq:temporal-strong}\\
 \Wtwo\big(\Law(X_{a,w}^{K}),\Law(X_{a,w}^{K,\delta})\big)
 &\le C\delta^\eta.
 \label{eq:temporal-w2}
\end{align}
For sufficiently small $\delta$, the corresponding CDF and quantile errors are bounded by $C\delta^\eta$ uniformly over the WP4 index sets.
\end{theorem}

\begin{proof}
By Cauchy--Schwarz,
\[
 |\Psi_{w-\Pi_\delta w}(P_KU)|^2
 \le
 \norm{w-\Pi_\delta w}_{\cH}^2
 \int_0^T\norm{P_KU_t}^2dt.
\]
The integrated moment is uniformly bounded by the baseline WP1 estimate. The coupling gives \eqref{eq:temporal-w2}, and Lemma~\ref{lem:normal-comparison} gives the law bounds after preserving the lower variance bound.
\end{proof}

\subsection{Left-point quadrature and approximate states}

Consider
\begin{equation}\label{eq:left}
 \bar X_{a,w}^{K,\delta}
 =\delta\sum_{j=0}^{J-1}
 \ip{w(t_j)}{U_{t_j}^{a,K}}.
\end{equation}

\begin{assumption}[Time regularity]\label{ass:time-holder}
For exponents $\eta_w,\chi>0$,
\begin{align}
 \sup_{w\in\cW}\norm{w_t-w_s}
 &\le C|t-s|^{\eta_w},
 &\sup_w\norm{w}_{L^\infty(0,T;H)}&<\infty,
 \label{eq:w-holder}\\
 \sup_{a,K}\E\norm{U_t^{a,K}-U_s^{a,K}}^2
 &\le C|t-s|^{2\chi}.
 \label{eq:U-holder}
\end{align}
A sufficient field condition for \eqref{eq:U-holder}, with $0<\chi\le1/2$, is
\begin{equation}\label{eq:holder-sufficient}
 u_0\in D(\Lam^\chi),
 \qquad
 \sum_kq_k^2\lambda_k^{2\chi-1}<\infty.
\end{equation}
\end{assumption}

\begin{proposition}[Robust rectangle-rule rate]\label{prop:left}
Under Assumption~\ref{ass:time-holder},
\begin{equation}\label{eq:left-rate}
 \sup_{a,K,w}
 \norm{X_{a,w}^{K}-\bar X_{a,w}^{K,\delta}}_{L^2}
 \le C\delta^{\min(\eta_w,\chi)}.
\end{equation}
If an approximate state $\widehat U^{K,\delta}$ satisfies
\begin{equation}\label{eq:state-scheme}
 \sup_{a,K}\E\int_0^T
 \norm{U_t^{a,K}-\widehat U_t^{a,K,\delta}}^2dt
 \le C\delta^{2\kappa_U},
\end{equation}
then the target-level rate is $\delta^\kappa$, where
\begin{equation}\label{eq:kappa}
 \kappa=\min\{\eta_w,\chi,\kappa_U\}.
\end{equation}
For the projected-weight scheme, replace $\eta_w$ and $\chi$ by the projection exponent $\eta$ and use $\kappa=\min\{\eta,\kappa_U\}$.
\end{proposition}

\begin{proof}
Let $\pi_\delta(t)$ be the left grid point. Then
\begin{align*}
 &\E\left|
 \ip{w_t}{U_t}-\ip{w_{\pi_\delta(t)}}{U_{\pi_\delta(t)}}
 \right|^2\\
 &\quad\le
 2\norm{w_t-w_{\pi_\delta(t)}}^2\E\norm{U_t}^2
 +2\norm{w_{\pi_\delta(t)}}^2
 \E\norm{U_t-U_{\pi_\delta(t)}}^2.
\end{align*}
Integrate in time and use $|\int_0^T f|^2\le T\int_0^T|f|^2$. The state-scheme term follows from Cauchy--Schwarz and \eqref{eq:state-scheme}.
\end{proof}

\begin{remark}[Scheme-specific burden]
The exact OU transition is exact at grid points. Implicit Euler, exponential Euler with approximate stochastic increments, and finite-element time steppers require their own stability and regularity estimates. Finite-dimensional convergence at fixed $K$ does not imply a $K$-uniform temporal rate. This is consistent with the classical time-discretisation analysis of parabolic SPDEs and with modern pathwise-uniform stochastic convolution estimates.
\end{remark}

\section{Bounded operator approximation}

Let the numerical generator be
\begin{equation}\label{eq:operator-approx}
 A_{a,M}=aA_0+E_{a,M},
 \qquad
 \sup_{a\in\widetilde\Theta}
orm{E_{a,M}}_{\cL(H)}\le\eta_M.
\end{equation}
This term covers bounded operator perturbations only. Spectral projection and domain-changing differential approximations remain governed by the Galerkin analysis and must not be hidden in $E_{a,M}$.

\begin{theorem}[Operator transfer and response correction]\label{thm:operator}
Let $T_A$ denote any non-degenerate Gaussian CDF, tail, or quantile target considered in WP6. Then
\begin{equation}\label{eq:op-first}
 \sup_\tau|T_{A_{a,M}}(\tau)-T_{aA_0}(\tau)|
 \le C\eta_M.
\end{equation}
If $E_{a,M}$ is known and
\begin{equation}\label{eq:op-corrected}
 T_{A_{a,M}}^{\mathrm{corr}}
 :=T_{A_{a,M}}-D_AT_{aA_0}[E_{a,M}],
\end{equation}
then
\begin{equation}\label{eq:op-second}
 \sup_\tau|T_{A_{a,M}}^{\mathrm{corr}}(\tau)-T_{aA_0}(\tau)|
 \le C\eta_M^2.
\end{equation}
For a scalar intervention grid $a_M$, take $\eta_M=|a_M-a|$ and use the WP6 analytic derivative $\partial_aT$.
\end{theorem}

\begin{proof}
Apply the WP6 uniform first-order operator Lipschitz theorem and its quadratic Fr\'echet remainder. The constants are uniform on the compact intervention and target classes.
\end{proof}

\section{Uniform primitive and loading approximation}

Let $T$ be any WP4 target family: fixed-intervention CDF, factual CDF, distributional effect, tail, or quantile surface. Set
\begin{equation}\label{eq:Bnum}
 B_{K,\delta,M}
 :=C\left(K^{-\rho_F}+\delta^\kappa+\eta_M\right).
\end{equation}
For a first-order operator-corrected implementation, replace $\eta_M$ by $\eta_M^2$.

\begin{theorem}[Target and derivative approximation]\label{thm:target-derivative}
Uniformly over the compact parameter rectangle,
\begin{align}
 \norm{T^{K,\delta,M}(\vartheta)-T(\vartheta)}_\infty
 &\le B_{K,\delta,M},
 \label{eq:target-num}\\
 \norm{DT^{K,\delta,M}(\vartheta)-DT(\vartheta)}_\infty
 &\le C\left(K^{-\rho_F}+\delta^\kappa+\eta_M\right).
 \label{eq:derivative-num}
\end{align}
The same spatial exponent applies to the scalar diffusivity derivatives through order two.
\end{theorem}

\begin{proof}[Proof structure]
The target bound is the triangle inequality applied to Theorems~\ref{thm:law-rate}, \ref{thm:temporal-projection} or Proposition~\ref{prop:left}, and Theorem~\ref{thm:operator}. For the intervention derivatives, the modal mean kernels contain $\lambda_k^jt^je^{-a\lambda_kt}$, $j=1,2$, whose $L^2(0,T)$ norm is $O(\lambda_k^{-1/2})$. The first and second derivative Volterra kernels have $L^1$ norm $O(\lambda_k^{-1})$, the same spectral order as the original kernel. Hence the tail proofs repeat with unchanged exponents. Temporal and bounded-operator derivative errors follow from the same continuous derivative maps used in WP4 and WP6.
\end{proof}

\section{Uniform total-error theorem}

The statistical normalisations are
\begin{equation}\label{eq:stat-rates}
 R_{N,n}^{\mathrm{fast}}=\bn:=\sqrt{Nn},
 \qquad
 R_{N,n}^{\mathrm{slow}}=\an:=\sqrt{\LamN},
 \qquad
 \LamN:=\sum_{k\le N}\lambda_k.
\end{equation}
The fast rate applies to fixed absolute-intervention and policy-to-policy targets. The slow rate applies to factual and factual-versus-counterfactual targets.

\begin{theorem}[Statistical--numerical decomposition]\label{thm:total}
Let $f$ be one of the WP4 target families and let $R_{N,n}^{(f)}$ be its rate in \eqref{eq:stat-rates}. Then
\begin{equation}\label{eq:total-bound}
 \norm{\widehat T_{N,n}^{K,\delta,M}-T_0}_\infty
 =\Op\big((R_{N,n}^{(f)})^{-1}\big)
 +O\left(K^{-\rho_F}+\delta^\kappa+\eta_M\right).
\end{equation}
Let
\begin{equation}\label{eq:bias-surface}
 b_{K,\delta,M}:=T_0^{K,\delta,M}-T_0.
\end{equation}
If \eqref{eq:derivative-num} tends to zero, then in $C(\cT_f)$,
\begin{equation}\label{eq:total-expansion}
 R_{N,n}^{(f)}
 \big(\widehat T_{N,n}^{K,\delta,M}-T_0\big)
 =\mathbb G_f
 +R_{N,n}^{(f)}b_{K,\delta,M}
 +\op(1),
\end{equation}
where $\mathbb G_f$ is the corresponding WP4 Gaussian process.
\end{theorem}

\begin{proof}
Decompose
\[
 \widehat T^{K,\delta,M}-T_0
 =\{T^{K,\delta,M}(\widehat\vartheta)-T^{K,\delta,M}(\vartheta_0)\}
 +b_{K,\delta,M}.
\]
The first term has the uniform WP4 Taylor expansion with the numerical derivative. Equation \eqref{eq:derivative-num} allows replacement of that derivative by the exact loading, and WP3 supplies the parameter limit. The second term is deterministic numerical bias. This proves \eqref{eq:total-expansion}; the norm bound follows immediately.
\end{proof}

\begin{corollary}[Numerical-bias trichotomy]\label{cor:trichotomy}
For a given family $f$:
\begin{enumerate}[label=(\roman*)]
 \item If $R_{N,n}^{(f)}\norm{b_{K,\delta,M}}_\infty\to0$, the centred WP4--WP5 limit is unchanged.
 \item If $R_{N,n}^{(f)}b_{K,\delta,M}\to b_f$ in $C(\cT_f)$, the limit is $\mathbb G_f+b_f$.
 \item If $R_{N,n}^{(f)}\norm{b_{K,\delta,M}}_\infty\to\infty$, centred root-information inference is invalid.
\end{enumerate}
\end{corollary}

\section{Root-information inference and Decision Gate 6}

\begin{theorem}[Family-specific undersmoothing]\label{thm:undersmoothing}
The WP5 confidence bands and tests remain valid for the full target if
\begin{align}
 \bn\left(K^{-\rho_F}+\delta^\kappa+\eta_M\right)&\to0
 &&\text{for fast targets},
 \label{eq:fast-under}\\
 \an\left(K^{-\rho_F}+\delta^\kappa+\eta_M\right)&\to0
 &&\text{for slow targets}.
 \label{eq:slow-under}
\end{align}
For an operator-response-corrected implementation, replace $\eta_M$ by $\eta_M^2$. A common implementation supporting both families must satisfy \eqref{eq:fast-under}.
\end{theorem}

\begin{theorem}[Bias-enlarged simultaneous bands]\label{thm:bias-bands}
Suppose a deterministic or high-probability certificate satisfies
\begin{equation}\label{eq:bias-cert}
 \norm{T_0^{K,\delta,M}-T_0}_\infty\le B_{K,\delta,M}.
\end{equation}
Add $B_{K,\delta,M}$ to every WP5 simultaneous half-width. The resulting band covers the full target with asymptotic probability at least $1-\alpha$, without requiring \eqref{eq:fast-under} or \eqref{eq:slow-under}.
\end{theorem}

\begin{proof}
On the event that the WP5 band contains the approximate target, the triangle inequality and \eqref{eq:bias-cert} imply that the enlarged band contains the full target. The first event has probability tending to $1-\alpha$.
\end{proof}

\begin{proposition}[Estimated bias correction]\label{prop:bias-correction}
If an estimated bias surface satisfies
\begin{equation}\label{eq:bias-est}
 R_{N,n}^{(f)}
 \norm{\widehat b_{K,\delta,M}-b_{K,\delta,M}}_\infty
 \overset{\Pp}{\longrightarrow}0,
\end{equation}
then
\begin{equation}\label{eq:bc-est}
 \widehat T^{\mathrm{bc}}
 :=\widehat T^{K,\delta,M}-\widehat b_{K,\delta,M}
\end{equation}
has the original centred WP4 limit and the WP5 bootstrap remains valid after applying the same correction to the target estimate.
\end{proposition}

\begin{remark}[Gate 6]
Decision Gate 6 is passed once either family-specific undersmoothing is imposed or a certified numerical-bias allowance is included in the simultaneous bands. The positive Gaussian core can therefore support full-field inferential claims after WP7.
\end{remark}

\section{Optimal numerical allocation}

Suppose
\begin{equation}\label{eq:error-model}
 \mathrm{Bias}(K,\delta,M)
 \lesssim
 C_sK^{-\rho_F}+C_t\delta^\kappa+C_oM^{-\rho_A},
\end{equation}
so that $\eta_M\asymp M^{-\rho_A}$.

\begin{proposition}[Balanced and undersmoothed choices]\label{prop:allocation}
For statistical normalisation $R$:
\begin{enumerate}[label=(\roman*)]
 \item Balancing each numerical term with $R^{-1}$ gives
 \begin{equation}\label{eq:balanced}
 K_{\mathrm{bal}}\asymp R^{1/\rho_F},
 \qquad
 \delta_{\mathrm{bal}}\asymp R^{-1/\kappa},
 \qquad
 M_{\mathrm{bal}}\asymp R^{1/\rho_A}.
 \end{equation}
 This is appropriate for first-order total-error control but generally produces local numerical bias at the inferential scale.
 \item For any slowly diverging $L_R$,
 \begin{equation}\label{eq:strict-under}
 K\asymp(RL_R)^{1/\rho_F},
 \qquad
 \delta\asymp(RL_R)^{-1/\kappa},
 \qquad
 M\asymp(RL_R)^{1/\rho_A}
 \end{equation}
 yields $R\,\mathrm{Bias}=O(L_R^{-1})$.
\end{enumerate}
\end{proposition}

\begin{proposition}[Cost-optimal bias shares]\label{prop:cost}
Suppose
\begin{equation}\label{eq:cost}
 \mathrm{Cost}(K,\delta,M)\asymp K^\zeta\delta^{-\xi}M^\upsilon.
\end{equation}
Allocate a total numerical-bias budget $\varepsilon$ through $x_s+x_t+x_o=\varepsilon$. Then
\[
 \mathrm{Cost}
 \asymp
 x_s^{-\zeta/\rho_F}
 x_t^{-\xi/\kappa}
 x_o^{-\upsilon/\rho_A}.
\]
The minimiser is
\begin{equation}\label{eq:shares}
 x_i=\frac{a_i}{a_s+a_t+a_o}\varepsilon,
 \qquad
 a_s=\frac\zeta{\rho_F},
 \quad
 a_t=\frac\xi\kappa,
 \quad
 a_o=\frac\upsilon{\rho_A}.
\end{equation}
\end{proposition}

\begin{proof}
Minimise the logarithm of the cost under the linear budget constraint. The first-order conditions give $a_i/x_i$ equal to a common Lagrange multiplier, which yields \eqref{eq:shares}.
\end{proof}

\section{Numerical validation}\label{sec:numerical}

The spatial experiment uses the one-dimensional heat spectrum $\lambda_k=\pi^2k^2$ and
\[
 r_0=0.8,
 \qquad
 \gamma=0.8,
 \qquad
 \beta=0.
\]
For constant-in-time modal weights, \eqref{eq:const-rate} and Corollary~\ref{cor:pointwise} predict
\[
 \rho_m^{\mathrm{smooth}}=2.6,
 \qquad
 \rho_V^{\mathrm{pt}}=4.6,
 \qquad
 \rho_{\mathrm{str}}=2.3.
\]
The fitted slopes are reported in Table~\ref{tab:slopes}.

\begin{table}[ht]
\centering
\small
\begin{tabular}{@{}lrr@{}}
\toprule
Quantity & Predicted slope & Fitted slope\\
\midrule
Deterministic mean tail & $-2.6$ & $-2.592$\\
Stochastic variance tail & $-4.6$ & $-4.627$\\
Strong functional RMS & $-2.3$ & $-2.316$\\
Optimal Gaussian $W_2$ & $-2.6$ & $-2.592$\\
Kolmogorov distance & $-2.6$ & $-2.592$\\
\bottomrule
\end{tabular}
\caption{Spatial-rate verification.}
\label{tab:slopes}
\end{table}

\begin{figure}[ht]
\centering
\includegraphics[width=.78\textwidth]{wp7_spatial_rates.png}
\caption{Spatial Galerkin errors. The strong shared-noise error is controlled by the square root of the variance tail, while optimal Gaussian law errors are controlled by the faster mean tail in this design.}
\label{fig:spatial}
\end{figure}

The temporal experiment uses $w_k(t)=c_kt$ and a piecewise-left projection of the deterministic weight while integrating every OU mode exactly. Mean, variance, optimal $W_2$, and Kolmogorov errors all display first-order convergence; fitted slopes range from $0.994$ to $0.999$.

\begin{figure}[ht]
\centering
\includegraphics[width=.78\textwidth]{wp7_temporal_rates.png}
\caption{First-order temporal law convergence for the exact-modal projected-weight scheme.}
\label{fig:temporal}
\end{figure}

A total-budget experiment uses $R$ as a generic information normalisation. Balanced choices $K\asymp R^{1/1.6}$, $\delta\asymp R^{-1}$, and $\eta_M\asymp R^{-1}$ keep $R$ times the CDF bias approximately constant at $1.5$. An undersmoothed design uses each numerical component at rate $R^{-1.7}$; its normalised bias falls from $0.184$ at $R=20$ to $0.0265$ at $R=320$.

\begin{figure}[ht]
\centering
\includegraphics[width=.78\textwidth]{wp7_total_error_balance.png}
\caption{Balanced approximation produces a local numerical bias; undersmoothing drives the normalised bias to zero.}
\label{fig:budget}
\end{figure}

The final experiment simulates the fast Gaussian scale estimator and constructs pointwise delta intervals for a fixed counterfactual CDF. Naive balanced intervals have severe undercoverage because numerical bias dominates the statistical half-width. Adding the certified exact bias gives conservative coverage near $0.975$. Under undersmoothing, naive coverage approaches the nominal level as the normalised bias vanishes.

\begin{figure}[ht]
\centering
\includegraphics[width=.78\textwidth]{wp7_bias_coverage.png}
\caption{Coverage implications of balanced bias, strict undersmoothing, and bias-enlarged intervals.}
\label{fig:coverage}
\end{figure}

\section{Completion statement and next package}

WP7 satisfies its completion gate:
\begin{enumerate}[label=(\roman*)]
 \item deterministic and stochastic Galerkin tails are explicit and uniform;
 \item strong and optimal Gaussian-law rates are separated;
 \item CDF and quantile rates are proved under the WP4 non-degeneracy condition;
 \item temporal approximation covers an exact-modal functional scheme and conventional left-point/state schemes;
 \item bounded operator error and the WP6 response correction enter the total bias;
 \item the fast and slow inferential normalisations are used correctly;
 \item undersmoothing, local-bias, dominant-bias, and bias-enlarged regimes are distinguished;
 \item cost-optimal numerical allocations are explicit;
 \item the predicted spatial and temporal slopes are reproduced numerically;
 \item Decision Gate 6 is passed for the positive Gaussian core.
\end{enumerate}

Seven of ten work packages are now complete. The remaining sequence is
\[
 \boxed{
 \mathrm{WP8}\longrightarrow\mathrm{WP9}\longrightarrow\mathrm{WP10}
 }.
\]
WP8 must prove the sharp non-identification boundary, particularly under unrestricted adapted forcing and under observationally equivalent factual mechanisms that induce distinct post-intervention laws.

\begin{thebibliography}{9}

\bibitem{Printems2001}
J. Printems.
\newblock On the discretization in time of parabolic stochastic partial differential equations.
\newblock \emph{ESAIM: Mathematical Modelling and Numerical Analysis}, 35(6):1055--1078, 2001.
\newblock doi:10.1051/m2an:2001148.

\bibitem{JentzenKloeden2009}
A. Jentzen and P. E. Kloeden.
\newblock Overcoming the order barrier in the numerical approximation of stochastic partial differential equations with additive space-time noise.
\newblock \emph{Proceedings of the Royal Society A}, 465(2102):649--667, 2009.
\newblock doi:10.1098/rspa.2008.0325.

\bibitem{vanNeervenVeraar2022}
J. van Neerven and M. Veraar.
\newblock Maximal inequalities for stochastic convolutions and pathwise uniform convergence of time discretisation schemes.
\newblock \emph{Stochastics and Partial Differential Equations: Analysis and Computations}, 10:516--581, 2022.
\newblock doi:10.1007/s40072-021-00204-y.

\bibitem{KovacsLarssonLindgren2012}
M. Kov\'acs, S. Larsson, and F. Lindgren.
\newblock Weak convergence of finite element approximations of linear stochastic evolution equations with additive noise.
\newblock \emph{BIT Numerical Mathematics}, 52:85--108, 2012.
\newblock doi:10.1007/s10543-011-0344-2.

\bibitem{DaPratoZabczyk2014}
G. Da Prato and J. Zabczyk.
\newblock \emph{Stochastic Equations in Infinite Dimensions}, second edition.
\newblock Cambridge University Press, 2014.

\end{thebibliography}

\end{document}
